\clearpage
\begin{proofbox}
\textbf{\textcolor{CProof}{证明之前先补一块：为什么右半边严格上升就够了？}}
\begin{center}\begin{tikzpicture}[font=\small,>=stealth,line join=round]
\draw[gray!40,rounded corners=3pt] (0,0) rectangle (14.2,8.3);\node at (7.1,7.700000000000001) {奇函数绕原点对称：左、中、右连成同向单调};
\begin{scope}[shift={(4.1,3.9499999999999997)},x=1.3636363636363635cm,y=0.32857142857142857cm]\draw[->,gray] (-2.2,0)--(2.2,0) node[right]{$x$};\draw[->,gray] (0,-7)--(0,7) node[above]{$y$};\draw[blue,thick,solid,domain=-1.8:1.8,samples=180] plot(\x,{(\x)^3});\fill[orange!85!black] (-1,-1) circle(2.2pt);\fill[orange!85!black] (0,0) circle(2.2pt);\fill[orange!85!black] (1,1) circle(2.2pt);\end{scope}\node[anchor=west,align=left,text width=5.0cm] at (8.4,4.0) {图取 $f(x)=x^3$\\[9pt]左侧 $f(u)<0$\\[9pt]原点 $f(0)=0$\\[9pt]右侧 $f(v)>0$\\[9pt]跨越 $0$ 仍有 $f(u)<f(v)$};
\node[text width=12.6cm,align=center] at (7.1,.65) {证明还要比较同在左侧的两点；用奇性将它们对应到右侧即可。};
\end{tikzpicture}\end{center}



设 $f(x)$ 为奇函数，在定义域右半边严格递增，且定义域是一个以 $0$ 为中间位置、左右对称的区间。
任取定义域内 $u<v$，讨论三种位置。
\begin{enumerate}[leftmargin=*,itemsep=0.3em]
  \item \textbf{$0\leqslant u<v$：}这就是题目已经给出的条件，$f(u)<f(v)$。
  \item \textbf{$u<v\leqslant0$：}因为 $-u>-v\geqslant0$，在右半边有 $f(-u)>f(-v)$。奇函数又告诉我们 $f(-u)=-f(u)$、$f(-v)=-f(v)$，所以 $-f(u)>-f(v)$，也就是 $f(u)<f(v)$。
  \item \textbf{$u<0<v$：}奇函数有 $f(0)=0$，严格递增又给出 $f(u)<0<f(v)$，当然有 $f(u)<f(v)$。
\end{enumerate}
所以奇函数在\textbf{整个定义域}严格递增。把所有“$<$”的方向反过来，同样说明：右半边严格递减的奇函数，在整个定义域也严格递减。

\medskip
\end{proofbox}\clearpage\begin{proofbox}\textbf{\textcolor{CProof}{证明一：最快的方法——搬到同一边比较（优先记）}}
\begin{center}\begin{tikzpicture}[font=\small,>=stealth,line join=round]
\draw[gray!40,rounded corners=3pt] (0,0) rectangle (14.2,8.3);\node at (7.1,7.700000000000001) {奇函数：把相加改成同一条曲线上的比较};
\begin{scope}[shift={(4.1,3.9499999999999997)},x=1.2cm,y=0.22999999999999998cm]\draw[->,gray] (-2.5,0)--(2.5,0) node[right]{$x$};\draw[->,gray] (0,-10)--(0,10) node[above]{$y$};\draw[blue,thick,solid,domain=-2.1:2.1,samples=180] plot(\x,{(\x)^3});\fill[orange!85!black] (1,1) circle(2.2pt);\fill[orange!85!black] (2,8) circle(2.2pt);\end{scope}\node[anchor=west,align=left,text width=5.0cm] at (8.4,4.0) {图取 $f(x)=x^3$\\[9pt]$a=2,\ b=-1$\\[9pt]$f(-b)=f(1)=1$\\[9pt]$f(a)+f(b)=8-1=7>0$\\[9pt]比较 $a=2>-b=1$};
\node[text width=12.6cm,align=center] at (7.1,.65) {图中标出的是比较用的 $a$ 和 $-b$；规律依赖奇性与严格单调性。};
\end{tikzpicture}\end{center}



\textbf{（1）奇函数，右半边严格递增：}
\[
\begin{aligned}
 f(a)+f(b)>0
 &\iff f(a)>-f(b) &&\text{先移项}\\
 &\iff f(a)>f(-b) &&\text{奇函数：$-f(b)=f(-b)$}\\
 &\iff a>-b &&\text{整个定义域严格递增}\\
 &\iff a+b>0. &&\text{移项}
\end{aligned}
\]
注意每一步都有双向箭头，说明不仅“由左能推出右”，而且“由右也能推回左”。如果只知道函数“不下降”，\textbf{第三步不能这样写}。

\textbf{（2）偶函数，右半边严格递增：}
\[
\begin{aligned}
 f(a)>f(b)
 &\iff f(|a|)>f(|b|) &&\text{偶函数：左右等高}\\
 &\iff |a|>|b|. &&\text{正半边严格递增}
\end{aligned}
\]

\textbf{（3）右半边严格递减：}奇函数在整个定义域严格递减，因此
\[
 f(a)+f(b)>0\iff f(a)>f(-b)\iff a<-b\iff a+b<0.
\]
偶函数同理：
\[
 f(a)>f(b)\iff f(|a|)>f(|b|)\iff |a|<|b|.
\]

\medskip

\end{proofbox}
\clearpage
\begin{proofbox}
\textbf{\textcolor{CProof}{证明二：分正负数讨论——初中生也能逐步核对}}
\begin{center}\begin{tikzpicture}[font=\small,>=stealth,line join=round]
\draw[gray!40,rounded corners=3pt] (0,0) rectangle (14.2,8.3);\node at (7.1,7.700000000000001) {奇函数绕原点对称：左、中、右连成同向单调};
\begin{scope}[shift={(4.1,3.9499999999999997)},x=1.3636363636363635cm,y=0.32857142857142857cm]\draw[->,gray] (-2.2,0)--(2.2,0) node[right]{$x$};\draw[->,gray] (0,-7)--(0,7) node[above]{$y$};\draw[blue,thick,solid,domain=-1.8:1.8,samples=180] plot(\x,{(\x)^3});\fill[orange!85!black] (-1,-1) circle(2.2pt);\fill[orange!85!black] (0,0) circle(2.2pt);\fill[orange!85!black] (1,1) circle(2.2pt);\end{scope}\node[anchor=west,align=left,text width=5.0cm] at (8.4,4.0) {图取 $f(x)=x^3$\\[9pt]左侧 $f(u)<0$\\[9pt]原点 $f(0)=0$\\[9pt]右侧 $f(v)>0$\\[9pt]跨越 $0$ 仍有 $f(u)<f(v)$};
\node[text width=12.6cm,align=center] at (7.1,.65) {证明还要比较同在左侧的两点；用奇性将它们对应到右侧即可。};
\end{tikzpicture}\end{center}



先证明奇函数严格递增的情况。因为 $f(0)=0$，故正数 $t$ 有 $f(t)>0$，负数 $t$ 有 $f(t)<0$。
\begin{enumerate}[leftmargin=*,itemsep=0.45em]
  \item \textbf{$a,b$ 都不小于 $0$。}此时 $f(a),f(b)$ 都不小于 $0$，且只要 $a,b$ 有一个不是 $0$，两者之和就大于 $0$。这与 $a+b>0$ 完全一致；两者都为 $0$ 时，两边都等于 $0$。
  \item \textbf{$a,b$ 都不大于 $0$。}同理 $f(a),f(b)$ 都不大于 $0$，它们的和与 $a+b$ 一样：若至少一个小于 $0$，两者都为负；若都为 $0$，两者都为 $0$。
  \item \textbf{一个非负、一个非正。}以 $a\geqslant0,b\leqslant0$ 为例：
  \[
   f(a)+f(b)=f(a)-f(-b).
  \]
  这里 $a$、$-b$ 都不小于 $0$，所以只需在右半边比较：
  \[
  \begin{array}{rcl}
    a>-b&\Longleftrightarrow&f(a)+f(b)>0,\\
    a=-b&\Longleftrightarrow&f(a)+f(b)=0,\\
    a<-b&\Longleftrightarrow&f(a)+f(b)<0.
  \end{array}
  \]
  而比较 $a$ 与 $-b$，实际上就是看 $a+b$ 的正负。若 $a\leqslant0,b\geqslant0$，交换两个字母即可。
\end{enumerate}
三类情况已经包括所有实数 $a,b$，因此结论成立。偶函数不必分正负：先将 $a,b$ 改写为 $|a|,|b|$，就在右半边比较。右半边严格递减时，以上比较方向全部反过来。

\medskip

\end{proofbox}
\clearpage
\begin{proofbox}
\textbf{\textcolor{CProof}{证明三：先找“恰好等于零”的分界点——适合解题}}
\begin{center}\begin{tikzpicture}[font=\small,>=stealth,line join=round]
\draw[gray!40,rounded corners=3pt] (0,0) rectangle (14.2,8.3);\node at (7.1,7.700000000000001) {固定一个输入：越过等号位置，和式才改变符号};
\begin{scope}[shift={(3.7341463414634153,3.2928571428571427)},x=1.4634146341463417cm,y=0.32857142857142857cm]\draw[->,gray] (-1.8,0)--(2.3,0) node[right]{$a$};\draw[->,gray] (0,-5)--(0,9) node[above]{$S$};\draw[blue,thick,solid,domain=-1.4:2,samples=180] plot(\x,{(\x)^3-1});\fill[orange!85!black] (1,0) circle(2.2pt);\end{scope}\node[anchor=west,align=left,text width=5.0cm] at (8.4,4.0) {取 $f(x)=x^3,\ b=-1$\\[9pt]纵轴为 $S(a)=f(a)+f(b)$\\[9pt]$S(a)=a^3-1$\\[9pt]$a<1$：负；$a>1$：正\\[9pt]分界点 $a=-b=1$};
\node[text width=12.6cm,align=center] at (7.1,.65) {横轴表示可变输入 $a$；图中和式的零点不是原函数的零点。};
\end{tikzpicture}\end{center}



\textbf{（1）奇函数：}先固定 $b$，只让 $a$ 改变。要让 $f(a)+f(b)=0$，等价于 $f(a)=f(-b)$。
由于函数严格递增或严格递减，同一个函数值只会对应一个自变量，所以\textbf{刚好相等的那个位置}就是
\[
 a=-b.
\]
若整个图像严格递增，$a$ 从 $-b$ 往右走，$f(a)$ 越来越大，于是 $f(a)+f(b)$ 变成正数；往左走则变成负数。若整个图像严格递减，正负恰好交换。

\textbf{（2）偶函数：}先固定 $b$，将 $|b|$ 看成 $0$ 右边的一个固定距离。当 $|a|=|b|$ 时，偶函数使两边函数值恰好相等；当 $|a|$ 比它大（或小），右半边严格递增保证函数值严格变大（或变小）。严格递减时反向。

\textbf{这个方法的优点：}先找等号成立的位置，再看越过这个位置后图像如何变化；即使具体的 $f(x)$ 不知道，也能判断不等号。

\medskip
\textbf{\textcolor{CProof}{额外补充：把“$>$”改成“$=$”或“$<$”}}

对于严格递增的奇函数，完整关系是
\[
\begin{array}{rcl}
 f(a)+f(b)>0&\iff&a+b>0,\\
 f(a)+f(b)=0&\iff&a+b=0,\\
 f(a)+f(b)<0&\iff&a+b<0.
\end{array}
\]
对于严格递增的偶函数，$f(a)$ 与 $f(b)$ 的大小关系，分别对应 $|a|$ 与 $|b|$ 的大小关系。严格递减时，\textbf{“大于”和“小于”互换，等于不变}。
\end{proofbox}
